\scanpage{101}
Je préfère faire les calculs du côté de $\mathcal A_1,\mathcal A_2$.
On peut expliciter la construction de $\mathcal A$ ainsi
\diagram{\begin{tikzcd}[column sep=large]
0 \arrow[r] & \overbrace{\mathcal A_1\unsure{\oplus|\amalg}_{\underline{\mathcal O}}\mathcal A_2}^{=\mathcal A_1\otimes_0\mathcal A_2} \arrow[r,no head] \arrow[d,"\tau"'] & \mathcal A_1\otimes\mathcal A_2 \arrow[r] \arrow[d,no head,"\unsure{\uparrow|\downarrow}"] & L_1^\vee\otimes L_2^\vee \arrow[r,no head] \arrow[d,equal] & 0\\
0 \arrow[r] & \underline{\mathcal O} \arrow[r] & \mathcal A \arrow[r] & L_1^\vee\otimes L_2^\vee \arrow[r] & 0
\end{tikzcd}}
\inserted{Trace réduite, qui est $\operatorname{Tr}_{A_1}$, $\operatorname{Tr}_{A_2}$ sur les deux composantes.}
$\mathcal A$ \inserted{$=\mathcal A_1*\mathcal A_2$} est l'ext. de $L_1^\vee\otimes L_2^\vee$ par $\underline{\mathcal O}$, déduite de l'extension $\mathcal A_1\otimes\mathcal A_2$ de $L_1^\vee\otimes L_2^\vee$ par $\mathcal A_1\otimes_0\mathcal A_2$, par l'hom $\tau$
\[
\mathcal A_1\amalg_{\underline{\mathcal O}}\mathcal A_2\mathrel{-\!\!-\!\!-}\mathcal O.
\]
Pour donner l'hom
\[
\mathcal A\longrightarrow\mathcal A_1\otimes\mathcal A_2
\]
c'est se donner séparément
\[
\underline{\mathcal O}\mathrel{-\!\!-\!\!-}\mathcal A_1\otimes\mathcal A_2
\]
\margin{(on prendra l'inclusion canonique $\lambda\mapsto\lambda.e_{A_1}\otimes e_{A_2}$}
et
\[
\mathcal A_1\otimes\mathcal A_2\xrightarrow[\inserted{x_1\otimes x_2\mapsto i(x_1*x_2)}]{\unsure{\varphi}}\mathcal A_1\otimes\mathcal A_2
\]
tel que les composés sur
\diagram{\begin{tikzcd}
& \mathcal O\\
\mathcal A_1\otimes_0\mathcal A_2=\mathcal A_1\oplus_{\mathcal O}\mathcal A_2 \arrow[ur] \arrow[dr] &\\
& \mathcal A_1\otimes\mathcal A_2
\end{tikzcd}}
soient égaux. On prendra
\mspage{100}
